\section{The Absurdity of Socialism}

\begin{quotationx}
To believe in the equality of all men, when we see them all unequal; to believe in liberty, when we see slavery established in all parts; to believe that all men are brothers, when history tells all are enemies; to believe that there is a common mass of misfortunes and of glories for all men born, when I see nothing but individual glories and misfortunes; to believe I am referred to humanity, when I know humanity is referred to me; to believe that humanity is my centre, when I constituted myself the centre of all; and finally, to believe that I should believe these things, when they are proposed to me by those who tell me that I should believe only my own reason, which contradicts all those things they propose to me, is an absurdity so stupendous, an abberation so inconceivable, that I stand mute and astounded in its presence.

My astonishment increases when I observe that those who affirm human solidarity, deny that of the family, which is to affirm that enemies are brothers, and that brothers should not be brothers; that those who affirm human solidarity are the same who a little before denied the political, which is to affirm I have nothing in common with my own, and all in common with strangers; that those who affirm human solidarity deny religion, though the former cannot be explained without the latter; and from all this I deduce in legitimate consequence that the Socialistic schools are at once illogical and absurd --- illogical, because after demonstrating against the Liberal school that some solidarities cannot be accepted while others are rejected, they fall into the same error, accepting one amongst all, and rejecting the remainder --- absurd, because precisely the one they proposed to me is not a point of reason but of faith, and because this proposal comes to me from those who deny faith and proclaim the imprescriptable right of reason to empire and sovereignty.

\flright{\textsc{Juan Donoso Cortés} (marqués de Valdegamas), \textit{Essays on Catholicism, Liberalism and Socialism}\footnote{\url{https://gornahoor.net/library/CortesEssays.pdf}}}
\end{quotationx}


\flrightit{Posted on 2009-08-29 by Cologero}

